\begin{table}[!htbp]
\centering
\caption{IV Estimates of Immigrant Inflows on U.S.-Born Internal Migration}
\label{tab:migration_iv}
\resizebox{\ifdim\width>\textwidth \textwidth\else \width\fi}{!}{

\begin{tabular}{lccc}
\toprule
  & All & Men & Women\\
\midrule
Levels (Card specification) & 6.96e-06 & 1.54e-06 & 3.03e-06\\
 & (4.46e-06) & (4.74e-06) & (3.73e-06)\\
Differences & 4.82e-05 & 8.61e-05 & 1.14e-05\\
 & (4.60e-05) & (6.11e-05) & (3.33e-05)\\
\midrule
Mean P(moved) & 0.1645 & 0.1678 & 0.1612\\
First-stage Wald, levels & 634.7 & 579.2 & 600.0\\
First-stage Wald, differences & 23.7 & 24.1 & 23.4\\
SE clustering & MSA & MSA & MSA\\
Observations, levels & 1,000,000 & 500,000 & 500,000\\
Observations, differences & 4,579,873 & 2,229,032 & 2,350,841\\
\bottomrule
\end{tabular}
}
\begin{minipage}{\textwidth}
\footnotesize
\vspace{4pt}
\textit{Notes:} The two specifications use different samples, hence the separate observation counts. The levels specification follows \citet{card2001immigrant} on a seeded random draw from U.S.-born respondents aged 16--68. The differences specification uses the full sample of U.S.-born respondents aged 25--68; its instrument differences a baseline window at ages 14--19 against the preceding five years, and those windows overlap below age 24, which would drive the instrument mechanically toward zero in a younger sample. Standard errors in parentheses are clustered on metropolitan area; state-clustered alternatives do not change any conclusion. The first-stage Wald statistic is computed against the same metropolitan-area-clustered covariance matrix as the standard errors rather than as a homoskedastic $F$: the instrument varies at the metropolitan-area-by-year level, so a statistic assuming independence across individual records would largely restate the number of records rather than measure instrument strength. $^{***}$ $p<0.01$, $^{**}$ $p<0.05$, $^{*}$ $p<0.10$.
\end{minipage}
\end{table}

